\documentclass{article}
\usepackage{amsmath,amssymb}
\usepackage{xurl}
\usepackage[hidelinks]{hyperref}
\setlength{\emergencystretch}{2em}
% Shared commands for notes in docs/paper-gaps/.

\DeclareUnicodeCharacter{2080}{\ifmmode{}_{0}\else\textsubscript{0}\fi}
\DeclareUnicodeCharacter{2081}{\ifmmode{}_{1}\else\textsubscript{1}\fi}
\DeclareUnicodeCharacter{2082}{\ifmmode{}_{2}\else\textsubscript{2}\fi}
\DeclareUnicodeCharacter{2099}{\ifmmode{}_{n}\else\textsubscript{n}\fi}

\newcommand{\Lean}{\textsc{Lean}}
\ExplSyntaxOn
\NewDocumentCommand{\leanid}{m}{
  \ifmmode
    \textsf{\detokenize{#1}}
  \else
    \group_begin:
    \urlstyle{sf}
    \tl_set:Nn \l_tmpa_tl {#1}
    \tl_replace_all:Nnn \l_tmpa_tl {\_} {_}
    \exp_args:NV \path \l_tmpa_tl
    \group_end:
  \fi
}
\ExplSyntaxOff
\providecommand{\tr}{\operatorname{tr}}
\newcommand{\tnleanIssueBase}{https://github.com/LionSR/TNLean/issues/}
\newcommand{\tnleanPrBase}{https://github.com/LionSR/TNLean/pull/}
\newcommand{\tnleanDocBase}{https://sirui-lu.com/TNLean/docs/}
\newcommand{\ghissue}[1]{\href{\tnleanIssueBase#1}{GitHub issue \##1}}
\newcommand{\ghpr}[1]{\href{\tnleanPrBase#1}{PR \##1}}
\newcommand{\doclink}[1]{\href{\tnleanDocBase#1.html}{\path{#1}}}

% Dirac notation and the identity operator, as in blueprint/src/macros/common.tex.
% \providecommand keeps note-local definitions authoritative.
\usepackage{dsfont}
\providecommand{\ket}[1]{|#1\rangle}
\providecommand{\bra}[1]{\langle #1|}
\providecommand{\braket}[2]{\langle #1 | #2 \rangle}
\providecommand{\ketbra}[2]{|#1\rangle\!\langle #2|}
\providecommand{\Id}{\mathds{1}}
\providecommand{\one}{\mathds{1}}
\providecommand{\C}{\mathbb{C}}
\providecommand{\MN}[1]{M_{#1}(\C)}
\providecommand{\MD}{M_D(\C)}

% Verdict marker: \gapnote{<kind>}{<status>}. Kind is the policy
% classification (clarification, local-correction, scope-restriction,
% unfaithful, false-source, open-gap); status is open, wip (elimination
% underway), resolved, or historical. Machine-read by the site index;
% prints a verdict line.
\newcommand{\gapnote}[2]{\par\noindent\textbf{Verdict:} #1 (#2).\par}

\providecommand{\leanid}[1]{\begingroup\urlstyle{sf}\path{#1}\endgroup}
\providecommand{\gapnote}[2]{\par\noindent\textbf{Verdict:} #1 (#2).\par}
\title{Constants in the AKLV Energy-Truncation Argument}
\author{}
\date{2026-10-29}
\begin{document}
\maketitle
\gapnote{local-correction}{resolved}

\section{Source and normalization}
Arad, Kitaev, Landau and Vazirani, arXiv:1301.1162v1,
Sections~6.1 and~6.3, consider positive nearest-neighbour terms bounded
by the identity. After regional scalar shifts the ground energy is in
$[0,6]$. Put $A=H_2+H_{s-1}$ and $B=H-A$.
In Section~6.3, $P_t=\mathbf 1_{(-\infty,t)}(B)$ has a strict threshold;
write $Q_t=1-P_t$. The truncation in Section~6.1 is the continuous
function $x\mapsto\min(x,t)$, so choosing either side at equality in its
spectral formula produces the same operator.

\section{The printed normalization denominator}
The uniform interior shift printed in Section~6.1 is $d/(s-7)$ for
$4\le i\le s-3$. For $s\ge8$ this range contains $s-6$ terms. If the
unshifted interior sum has ground energy $d>0$, the printed shifts give
ground energy
\begin{align}
    d-(s-6)\frac{d}{s-7}=-\frac{d}{s-7}<0,
\end{align}
rather than the stated zero. To obtain zero, the uniform shift would
be $d/(s-6)$; equivalently, subtract $dI$ from the interior sum once.
This correction follows by counting the displayed range. The normalized
robustness theorem takes the regional positivity and $0\le e\le6$ as
data and does not use the printed distribution of scalar shifts.

\section{The cross term and two typographical slips}
The Cauchy--Schwarz calculation in the proof of Lemma~6.5 gives
\begin{align}
    |\langle P_tv,AQ_tv\rangle|
    &\le \|A\|\,\|P_tv\|\,\|Q_tv\|
    \le 2\sqrt{\langle v,H^{(t)}v\rangle/t}
\end{align}
for a unit vector and positive $t$, using the first assertion of that lemma.
The last expression bounds the unsquared modulus in this calculation.
Squaring gives the different estimate
$|\langle P_tv,AQ_tv\rangle|^2\le4\langle v,H^{(t)}v\rangle/t$.
The unsquared estimate is proved by
\leanid{TNLean.MPS.AreaLaw.AKLV.norm_low_high_boundary_inner_le_sqrt_clipped_energy}.
The actual mass estimate is proved from positivity and pairwise commutation
of the three disconnected regional blocks by
\leanid{TNLean.MPS.AreaLaw.AKLV.norm_highProjection_le_sqrt_clipped_energy}.
The proof of Lemma~6.6(1) uses $\epsilon_0$, where the printed expansion
has $\epsilon$. Positivity of $A$ gives
$\langle Q_t\Gamma,HQ_t\Gamma\rangle\ge
\langle Q_t\Gamma,BQ_t\Gamma\rangle$; the comparison printed in the prose
after that expansion has the opposite sign.

\section{Recovering the off-diagonal constant}
The printed commutator bound $\|Q_i\|\le2\cdot4^i i!$ sums at $r=1/8$
to $4$, while Lemma~6.6(2) states coefficient $2$.
For each boundary bond replace $h_j$ by $O_j=h_j-I/2$.
Then $\|O_j\|\le1/2$, and the local Hadamard bound gives
$\|e^{B/8}O_je^{-B/8}\|\le1$.
For $t\ge u$, $Q_tP_u=0$, so the scalar shifts disappear and
\begin{align}
    \|Q_tAP_u\|&\le e^{-(t-u)/8}
      \sum_{j\in\{2,s-1\}}\|e^{B/8}O_je^{-B/8}\|
      \le2e^{-(t-u)/8}.
\end{align}
The proved statement is
\leanid{TNLean.MPS.AreaLaw.AKLV.norm_two_boundary_spectral_offDiagonal_le}.
Its hypotheses are the original local positive contractions and their
supports; no commutator-growth or spectral-tail bound is supplied.
The complement may have an arbitrary real scalar shift.

\section{An explicit eigenvector tail}
Let $(B+A)x=ex$, $0\le e\le6$, $W=e^{B/8}$ and $y=Wx$.
The uncentered local estimate gives $\|WAW^{-1}\|\le4$, hence
$\|By\|\le10\|y\|$. The spectral inverse above energy $20$ gives
$\|Q_{20}y\|\le\|By\|/20\le\|y\|/2$.
The low part satisfies $\|P_{20}y\|\le e^{20/8}\|x\|$, so
$\|y\|\le2e^{20/8}\|x\|$. For every real $t$,
\begin{align}
    \|Q_tx\|&\le\|Q_te^{-B/8}\|\,\|y\|
      \le2e^{-(t-20)/8}\|x\|.
\end{align}
This alternative to the shell induction of Lemma~6.7 is proved by
\leanid{TNLean.MPS.AreaLaw.AKLV.norm_groundState_highProjection_le_exp}.
It applies to every eigenvector in this energy range, including the unit
ground vector. It requires neither a gap nor uniqueness.

\section{Explicit robustness constants and the final statement}
Separate clipping of the two exterior blocks, its positive order bounds,
and its agreement with the original operator on the low-energy subspace
use QICLean's energy truncation. The full Hamiltonian is not clipped in
place of the exterior blocks, and no ALVV regional approximation theorem
is substituted for this construction.

The projected-energy comparison is
$\langle P_tv,HP_tv\rangle\le\langle v,H^{(t)}v\rangle+4\|Q_tv\|$.
For $tg^2\ge1024(e+g)$ it implies that every unit vector orthogonal to
the original ground vector has clipped energy at least $e+g/2$.
A min--max argument constructs the actual clipped ground vector and
proves its gap is at least $g/2$.

The projected original ground vector has residual at most
$10\exp(-(t-40)/16)$. Combining this with the absolute excited-energy
threshold gives a phase-aligned squared-distance bound
$968g^{-2}\exp(-(t-40)/8)$.
For $t\ge32768(e+g)/g^2$, this is at most $\exp(-t/16)$.
The final statement is
\leanid{TNLean.MPS.AreaLaw.AKLV.normalized_chain_robustness}.

Its hypotheses are the original nearest-neighbour supports and
$0\le h_i\le I$, the two boundary bonds, the three scalar energy shifts,
positivity of the resulting regional blocks, a unit original ground
vector with $0\le e\le6$, and the original gap $0<g\le1$ in operator
form. The three positivity assumptions and energy bound are the
normalization fixed in Section~6.1. Regional commutations and the sum
decomposition are derived from the actual supports and cuts. No
spectral tail, commutator-growth bound, clipped gap, residual estimate
or closeness of the ground vectors is supplied as an assumption.
The conclusion uses the actual separate clipping, an actual ground
eigenvector, the operator gap inequality, and an explicit complex phase.
Constructing the normalization from an unshifted Hamiltonian is the
source's preceding reduction; this theorem is stated on the normalized
data of Theorem~6.1.
\end{document}
